%% notes-tex/027-metabolic-flux-module/sections/1-what-it-is.tex
%% [[experiments.027-betaxanthin-metabolic]]
%% https://github.com/Mjvolk3/torchcell/tree/main/notes-tex/027-metabolic-flux-module/sections/1-what-it-is.tex

\section{What the module computes\secstatus{todo}}
\label{sec:what}

\subsection{Where it sits\secstatus{todo}}

The Equivariant Cell Graph Transformer separates its instruments by what they map between.
A \textbf{Type~I} instrument is representation to representation: the perturbation
operators $T_\psi^{\mathrm{del}}$, $T_\psi^{\mathrm{OE}}$, $T_\psi^{\mathrm{KD}}$ act
inside the gene embedding space and hand on $H_{\mathrm{pert}}$. A \textbf{Type~II}
instrument is representation to output: a readout that emits a measured phenotype.

The flux module is neither, and the reason is worth stating precisely rather than settling
by convention. By the letter of the taxonomy it is Type~I, since its output feeds further
computation rather than a loss. But every existing Type~I instrument acts \emph{within} the
gene index space at dimension $d$ and preserves equivariance over genes, whereas this one
leaves gene index space entirely for reaction space $\mathbb{R}^{4131}$ under a fixed,
externally supplied stoichiometry, and its parameters are tabulated physical constants
rather than learned weights. It is a physically identified decoder, and its output
$\nu$ is a latent physical state rather than a phenotype: the phenotypes are read
\emph{off} it.

\subsection{A program cannot learn, which is the point\secstatus{todo}}

The constraint-based tradition solves a linear program per condition. A linear program has
no use for a phenotype measurement, so it cannot be improved by one. This layer inverts
that: every constraint is either enforced exactly by the parameterization or relaxed to a
smooth penalty, so the flux is \emph{conditioned} on roughly $10^7$ fitness records, 4,735
betaxanthin measurements and 4,678 amino-acid profiles. \textbf{Nothing is enforced by a
binary variable}, which is what keeps the whole object differentiable and lets a batch of
perturbations cost the same forward pass as a single one.

\subsection{The relaxation, term by term\secstatus{todo}}

\begin{table}[H]
  \centering
  \footnotesize
  \caption[The relaxation]{Each constraint of an enzyme-constrained thermodynamic model,
  and how this layer realizes it. ``Exact'' means the parameterization cannot violate it;
  ``soft'' means a penalty that the optimizer trades against the data term.}
  \label{tab:relaxation}
  %% SOURCE: torchcell/metabolism/flux_layer.py, module docstring table and forward().
  \begin{tabular}{@{}p{0.26\textwidth}p{0.34\textwidth}p{0.32\textwidth}@{}}
    \hline
    Constraint & Conventional form & Here \\
    \hline
    box and directionality & $lb_j \le \nu_j \le ub_j$ & \textbf{exact}, via sigmoid \\
    enzyme capacity & $\nu_j \le k_{cat} E$ & \textbf{exact}, folded into the box \\
    mass balance & $S\nu = 0$ & soft, scale-relative \\
    second law & $\Delta_r G_j \nu_j < 0$ plus a binary & soft hinge, no binary \\
    Gibbs dissipation & $g_{\mathrm{diss}} \le g_{\mathrm{lim}}$ & soft hinge \\
    protein budget & $\sum_g MW_g E_g \le P_{\mathrm{avail}}$ & soft hinge \\
    \hline
  \end{tabular}
\end{table}

\paragraph{From representation to gene availability.} Each gene's embedding passes through
a learned gate, $\gamma_g = \sigma(w^\top h_g)$, so a gene is not merely present or absent.
A deletion pins $\gamma_g = 0$; everything else is free to take an intermediate value,
which is what leaves room for dosage, hypomorphic alleles and overexpression rather than
forcing the layer to treat every perturbation as a knockout.

\paragraph{Folding up the gene-protein-reaction rule.} Availability is folded to the
reaction through the GPR: a \textbf{softmin} over the subunits of a complex, since a
complex is limited by its scarcest member, and a \textbf{sum} over isozymes, since they
substitute for one another. The softmin temperature is a real parameter, not a detail: too
soft and a complex behaves like a mean over its subunits, which loses exactly the
limiting-subunit behavior the softmin was chosen for.

\paragraph{The box, which is where capacity lives.} Flux is parameterized as
\begin{equation}
  \nu_j = lb_j + (ub_j - lb_j)\,\sigma(z_j),
  \qquad
  |\nu_j| \le k_{cat,j}\,E_{g(j)},
\end{equation}
so the bounds and the capacity bound hold by construction and cannot be violated no matter
what the network does. Turnover enters in units of $\mathrm{h}^{-1}$
($k_{cat}\,[\mathrm{s}^{-1}] \times 3600$), which is what fixes the units of the enzyme
abundance $E$. The exactness budget is spent here rather than on $S\nu = 0$, and that is a
deliberate choice: bounds are exact and mass balance is soft, not the reverse.

\paragraph{The second law, which is the load-bearing relaxation.}
\begin{equation}
  C_{\mathrm{th}} = \frac{1}{|\mathcal{J}|}\sum_{j\in\mathcal{J}}
    \operatorname{relu}\!\big(\nu_j\,\Delta_r G_j + \epsilon\big),
\end{equation}
with $\mathcal{J}$ excluding exchange reactions, the biomass reaction and water transport,
which cross the system boundary or are lumped and so carry no meaningful reaction energy.
The $\epsilon$ is not decoration. The conventional big-M pair admits
$\nu_j = \Delta_r G_j = 0$ under either value of its binary, so the strict inequality is
never actually realized, and a smooth version has to name its tolerance rather than inherit
a fiction.

\textbf{Loop-freedom comes free, and that is the trick.} Because $\Delta_r G$ is a
difference of potentials it sums to zero around any cycle, so no cycle can run downhill at
every step. Requiring $\nu_j \Delta_r G_j \le 0$ therefore forbids internal loops with no
integer variable at all, which is precisely the property that makes loopless flux balance
analysis a mixed-integer program.

\subsection{Free against anchored thermodynamics\secstatus{todo}}

The in-cell reaction energy decomposes into three terms,
\begin{equation}
  \Delta_r G_j
  = \underbrace{\Delta_r G'^{\circ}_j}_{\text{tabulated}}
  + \underbrace{RT\!\!\sum_{i \neq \mathrm{H}^+}\!\! S_{ij} \ln c_i}_{\text{concentration}}
  + \underbrace{\Delta_r G^{t}_j}_{\text{transport}} .
\end{equation}

The distinction between the two thermodynamic arms is exactly which of these is supplied.
In the \textbf{free} arm the potential is learned, so the layer gets loop-freedom and
nothing else; the constraint is structural. In the \textbf{anchored} arm the standard term
is read from the tabulated formation energies, so the sign of every reaction is pinned by
measured chemistry rather than by whatever the optimizer finds convenient. That difference
is the experiment: structural loop-freedom is a real constraint on its own, and the
question is whether the chemistry adds anything to it.
